Ce travail pratique a permis de caractériser un prisme en verre par la méthode du minimum de déviation. Les mesures expérimentales ont permis d'établir les valeurs suivantes :
\begin{itemize}
    \item Angle du prisme : $A = 60,3^\circ \pm 0,5^\circ$
    \item Angle de déviation minimale : $D_m = 60,5^\circ \pm 0,5^\circ$
\end{itemize}

L'indice de réfraction calculé pour la raie jaune du sodium ($\lambda = 589,3$ nm) est de $n \approx 1,733$. Cette valeur est cohérente avec la nature du composant (verre de type Flint) 
et valide la précision du protocole malgré une lecture limitée au degré sur le goniomètre. 